% Kapitel 6: Fazit und Ausblick
\section{Fazit und Ausblick}
\label{sec:fazit}

\subsection{Beantwortung der Leitfrage}

Die Ausgangsfrage dieser Arbeit lautete: \textit{Wie beeinflusst die Netzwerkstruktur eines Smart Grids die Fehlertoleranz beziehungsweise die Robustheit gegenüber zufälligen Ausfällen und gezielten Angriffen?}

Die Ergebnisse lassen sich in drei Punkten zusammenfassen.

Erstens hat die Topologie des Netzes einen erheblichen Einfluss auf die Robustheit. Homogene Netze (ER) sind gegenüber zufälligen Ausfällen deutlich widerstandsfähiger als skalenfreie Netze (BA). Bei gezielten Angriffen verstärkt sich dieser Effekt noch weiter, da die Hub-Struktur skalenfreier Netze einen klaren Angriffspunkt bietet.

Zweitens macht die Art der Störung einen fundamentalen Unterschied. Die AUC-Differenz zwischen zufälligen und gezielten Angriffen auf ein BA-Netz beträgt über 0,5. Für die Praxis bedeutet das, dass die Robustheit gegenüber zufälligen Ausfällen keine Rückschlüsse auf die Widerstandsfähigkeit gegenüber gezielten Angriffen erlaubt.

Drittens zeigt sich, dass nicht die Interdependenz an sich das Kernproblem darstellt, sondern deren Struktur. Ein Smart Grid mit zufällig zugewiesenen Abhängigkeiten bleibt selbst bei vollständiger Kopplung relativ robust. Erst wenn die Abhängigkeiten so strukturiert sind, dass wichtige Stromknoten von wichtigen Kommunikationsknoten abhängen, wird ein Angriff auf das Kommunikationsnetz genauso wirkungsvoll wie ein direkter Angriff auf das Stromnetz.

\subsection{Weitergehende Erkenntnisse}

Neben der Beantwortung der Leitfrage ergeben sich einige zusätzliche Erkenntnisse.

Redundanz bei den Abhängigkeitsbeziehungen kann die Robustheit substanziell erhöhen. Bereits eine Verdopplung der Abhängigkeiten ($r = 2$) reicht aus, um die AUC nahezu auf das Niveau eines ungekoppelten Netzes zu bringen. Das deutet darauf hin, dass auch vergleichsweise einfache Schutzmaßnahmen einen spürbaren Effekt haben können.

Außerdem beeinflusst die Wahl des Netzwerkmodells die Simulationsergebnisse stark. Analysen, die ausschließlich auf homogenen ER-Netzen basieren, unterschätzen die Verwundbarkeit realer Systeme möglicherweise erheblich. Für aussagekräftige Robustheitsstudien ist es daher empfehlenswert, verschiedene Topologien einzubeziehen, insbesondere solche mit heterogenen Gradverteilungen.

\subsection{Ausblick}

Mehrere Aspekte konnten im Rahmen dieser Arbeit nicht behandelt werden und bieten Ansatzpunkte für weiterführende Untersuchungen.

Ein naheliegender Schritt wäre die Einführung kaskadierender Ausfälle, bei denen der Ausfall von Stromknoten rückwirkend auch Kommunikationsknoten beeinträchtigt. Das würde die Dynamik interdependenter Netze realistischer abbilden und die beobachteten Effekte möglicherweise noch verstärken.

Ein weiterer Ansatzpunkt wäre die Anwendung des Modells auf reale Netztopologien, etwa auf das europäische Hochspannungsnetz oder auf regionale Verteilnetze, für die öffentlich zugängliche Daten existieren. Damit ließe sich überprüfen, inwieweit die in dieser Arbeit identifizierten Mechanismen auch unter realistischeren Bedingungen Bestand haben.

Schließlich bleibt die Frage offen, wie sich die korrelierte Abhängigkeitszuweisung in Kombination mit Redundanz verhält. Die Ergebnisse zeigen, dass Redundanz bei zufälliger Zuweisung gut schützt. Ob das in gleichem Maße auch bei korrelierter Struktur gilt, wo beide redundanten Abhängigkeiten eines Stromknotens mit hoher Wahrscheinlichkeit ebenfalls an Hub-Kommunikationsknoten geknüpft sind, wäre eine lohnende Fragestellung für zukünftige Arbeiten.
